\section{Directional widths determine the full spectrum}
\label{sec:widths}

The diameter law uses only the largest and smallest scales of a
sublevel set. The intermediate eigenvalues depend on how many independent
directions have each scale. The radial widths of the difference body
record this information without reference to a barrier or a center.

For $0<g\leq\Delta$ and a Euclidean unit vector $h\in\V$, define the radial function
of the difference body by
\begin{equation}
 q_g(h)=\max\{t\geq0:th\in K_g\}.
 \label{eq:radial-width}
\end{equation}
The set $K_g$ is compact, convex, symmetric, and contains a neighborhood
of zero. Thus $q_g$ is positive, finite, even, and continuous on the unit
sphere. Continuity follows because $q_g$ is the
reciprocal of the Minkowski functional of $K_g$, a norm equivalent to
the Euclidean norm. For $j=1,\ldots,d$, define
the radial widths
\begin{align}
 w_j^-(g)&=\inf_{\substack{U\subset\V\\\operatorname{codim}U=j-1}}
       \ \sup_{\substack{h\in U\\\norm h=1}}q_g(h),
       \label{eq:radial-width-minus}\\
 w_j^+(g)&=\sup_{\substack{S\subset\V\\\dim S=j}}
       \ \inf_{\substack{h\in S\\\norm h=1}}q_g(h).
       \label{eq:radial-width-plus}
\end{align}
Here and below $U,S$ are linear subspaces. These two radial widths
need not coincide for a general convex body.

\begin{theorem}[Radial width bounds]
\label{thm:radial-spectrum}
Under the hypotheses of \cref{thm:difference-body}, list the
eigenvalues of $H_F(x)$ in
nondecreasing order. Then for $j=1,\ldots,d$,
\begin{equation}
 \frac1{(w_j^+(g))^2}\leq\lambda_j(H_F(x))
       \leq\frac{4C^2}{(w_j^-(g))^2}.
 \label{eq:radial-spectrum}
\end{equation}
Moreover, $w_j^+(g)\leq w_j^-(g)\leq2Cw_j^+(g)$,
$w_1^+(g)=w_1^-(g)=D(g)$, and
$w_d^+(g)=w_d^-(g)=r_g$.
\end{theorem}

\begin{proof}
The radial form of \eqref{eq:difference-body-sandwich} is
\[
 \frac1{\sqrt{\inner{H_F(x)h}{h}}}
 \leq q_g(h)\leq
 \frac{2C}{\sqrt{\inner{H_F(x)h}{h}}},\qquad \norm h=1.
\]
Equivalently, for the Rayleigh quotient $R(h)=\inner{H_F(x)h}{h}$,
\begin{equation}
 q_g(h)^{-2}\leq R(h)\leq4C^2q_g(h)^{-2}.
 \label{eq:radial-rayleigh}
\end{equation}
The two Courant--Fischer formulas in increasing eigenvalue order are
\[
 \lambda_j=\sup_{\operatorname{codim}U=j-1}
                   \inf_{h\in U,\,\norm h=1}R(h),\qquad
 \lambda_j=\inf_{\dim S=j}
                   \sup_{h\in S,\,\norm h=1}R(h).
\]
Apply the lower estimate in the second formula and the upper estimate
in the first. Positivity and uniform boundedness away from zero justify
taking reciprocal squared infima and suprema. Explicitly,
\begin{align*}
 \lambda_j&\geq\inf_{\dim S=j}\sup_{h\in S,\,\norm h=1}q_g(h)^{-2}
   =\left(\sup_{\dim S=j}\inf_{h\in S,\,\norm h=1}q_g(h)\right)^{-2},\\
 \lambda_j&\leq4C^2\sup_{\operatorname{codim}U=j-1}
                    \inf_{h\in U,\,\norm h=1}q_g(h)^{-2}
   =4C^2\left(\inf_{\operatorname{codim}U=j-1}
                    \sup_{h\in U,\,\norm h=1}q_g(h)\right)^{-2}.
\end{align*}
These are \eqref{eq:radial-spectrum}.

Every $j$-dimensional $S$ intersects every codimension-$(j-1)$ subspace
$U$ in a nonzero direction. A unit vector in that intersection shows
$\inf_{h\in S}q_g(h)\leq\sup_{h\in U}q_g(h)$; taking the indicated
extrema proves $w_j^+\leq w_j^-$. Combining the two bounds in
\eqref{eq:radial-spectrum} gives $w_j^-\leq2Cw_j^+$.
For $j=1$ both widths are the maximum
radial extent, namely $D(g)$. For $j=d$ both are the minimum radial
extent, namely the centered inradius $r_g$.
\end{proof}

For an ellipsoid, both radial widths are its principal semiaxes in
decreasing order, and the theorem reduces to the reciprocal-square
relation between semiaxes and Hessian eigenvalues. For a general
difference body, \eqref{eq:radial-spectrum} is the corresponding
variational statement. It can separate several spectral scales, which a
single condition number cannot.

\subsection{Exit distances along rays and a sharper constant at a given point}

A second pair of widths, measured from $x$ to the relative boundary
along one ray per line direction, improves the factor $4C^2$ to $C^2$.
Fix
$x\in P^\circ$. For a line $[h]\subset\V$ through the origin, choose
its Euclidean unit representative $\widehat h$ so that $\inner{c}{\widehat h}<0$ whenever
the objective is nonconstant on that line. If the line is orthogonal to
$c_{\V}$, choose the sign for which the ray from $x$ stays longer in
$P^\circ$, breaking ties arbitrarily. Define
\begin{equation}
 \tau_x([h])=\sup\{t\geq0:x+s\widehat h\in P^\circ
                                  \text{ for }0\leq s<t\}.
 \label{eq:directed-exit}
\end{equation}
The chosen ray stays in $L(g(x))$ until it leaves $P^\circ$, so its exit
point lies in $L(g(x))$, and $0<\tau_x([h])\leq\ell(x)$; this also
holds in the unbounded setting of \cref{prop:localization}. The sign
convention can make $\tau_x$ discontinuous at lines orthogonal to the
objective; this is why the next definitions use infima and suprema. Put
\begin{align}
 v_j^-(x)&=\inf_{\operatorname{codim}U=j-1}
            \sup_{h\in U,\,\norm h=1}\tau_x([h]),\label{eq:exit-width-minus}\\
 v_j^+(x)&=\sup_{\dim S=j}
            \inf_{h\in S,\,\norm h=1}\tau_x([h]).\label{eq:exit-width-plus}
\end{align}

\begin{corollary}[Exit-distance spectral bounds]
\label{cor:exit-spectrum}
Under the hypotheses of \cref{thm:difference-body},
\begin{equation}
 \frac1{(v_j^+(x))^2}\leq\lambda_j(H_F(x))
         \leq\frac{C^2}{(v_j^-(x))^2},\qquad j=1,\ldots,d.
 \label{eq:exit-spectrum}
\end{equation}
Also $v_j^+(x)\leq v_j^-(x)\leq Cv_j^+(x)$ and
$v_1^+(x)=v_1^-(x)=\ell(x)$.
\end{corollary}

\begin{proof}
The exit point lies at local-norm distance at least one from $x$ by
open Dikin containment, and at most $C$ by
\eqref{eq:sublevel-containment}.
Consequently, for every unit $h$,
\[
 \tau_x([h])^{-2}\leq\inner{H_F(x)h}{h}
                      \leq C^2\tau_x([h])^{-2}.
\]
All exit times are bounded below by the radius of a fixed relative ball
about $x$ and above by $\ell(x)$. The minimax argument in the proof of
\cref{thm:radial-spectrum} therefore applies, without continuity, and
also proves $v_j^+\leq v_j^-$. Combining the two bounds in
\eqref{eq:exit-spectrum} gives $v_j^-\leq Cv_j^+$.
Finally, for any $y\in L(g(x))$ distinct
from $x$, the ray along $y-x$ reaches at least as far as $y$. The selected
sign agrees with that ray when the objective strictly decreases, and
the longer-exit convention can only increase its length when the
objective is constant. Thus
$\sup_{[h]}\tau_x([h])\geq\norm{y-x}$. Maximizing over $y$ and using
the reverse inequality already noted proves
$v_1^+=v_1^-=\ell(x)$.
\end{proof}

The radial widths depend only on the sublevel geometry, so they allow
comparison across barriers. The exit-distance bounds depend on the
iterate and give sharper constants there. Neither
formulation assumes that the optimal face is nondegenerate or that the
spectrum has only two clusters.
